\chapter{Vaping}
\index{Vaping}

\section*{Definition and Overview}
Vaping is the act of inhaling and exhaling an aerosol (often incorrectly referred to as vapour) produced by an electronic nicotine delivery system (ENDS), electronic non-nicotine delivery system (ENNDS), or other vaporizing device. These devices, commonly known as e-cigarettes, vapes, vape pens, mods, or pod systems, operate by heating a liquid solution (e-liquid or vape juice) to temperatures sufficient to create an aerosol. Unlike traditional combustible tobacco products, which burn tobacco leaf to produce smoke, vaping devices rely on electrical resistance heating elements powered by lithium-ion batteries. The e-liquid typically consists of a solvent carrier (most commonly vegetable glycerin and propylene glycol), synthetic or tobacco-derived nicotine, and a complex mixture of flavoring chemical agents. Some devices are also modified or specifically designed to vaporize tetrahydrocannabinol (THC), synthetic cannabinoids, or various herbal extracts.

Initially introduced to global markets in the mid-2000s as a smoking cessation aid and a presumably safer alternative to traditional smoking, vaping has transitioned into a major public health challenge. The rapid technological evolution of these devices---moving from early "cig-a-like" models to high-capacity, rechargeable tank systems and discreet, high-nicotine pod-based devices---has drastically increased their appeal and accessibility. Vaping is highly significant from a clinical standpoint due to its potential to induce severe nicotine addiction, its direct toxic effects on the respiratory and cardiovascular systems, and its role as a gateway to combustible tobacco use among youth. Furthermore, the emergence of acute clinical syndromes, such as e-cigarette or vaping product use-associated lung injury (EVALI), has underscored the acute biological risks associated with the inhalation of vaporized chemical mixtures, challenging the initial narrative of these devices as harmless consumer products.

\section*{Epidemiology}
The epidemiological landscape of vaping has shifted dramatically over the past two decades, characterized by an exponential increase in prevalence, particularly among adolescents and young adults. Globally, millions of individuals engage in vaping, with the highest concentration of users in North America, Western Europe, and Australasia. According to national health surveys, vaping rates among high school students in some jurisdictions peaked at alarmingly high levels in the late 2010s and early 2020s, with up to 20\% to 30\% of adolescents reporting past-month use. Although regulatory interventions and public health campaigns have led to modest declines in some demographics, e-cigarette use remains the most common form of tobacco product consumption among youth in many developed countries.

Demographically, vaping exhibits distinct patterns. While traditional cigarette smoking is historically more prevalent among older populations and lower socioeconomic groups, vaping has captured a younger, more diverse demographic. The highest prevalence is consistently observed in the 15--24 age cohort. There is also a significant gender disparity in certain regions, with young males historically reporting higher rates of daily use, though this gap has narrowed due to the aggressive marketing of flavored, aesthetically appealing pod vapes targeted at all youth. Socioeconomic status, educational attainment, and geographic location also influence usage patterns, with higher rates often observed in suburban and urban areas where retail access is concentrated. Furthermore, there is a strong epidemiological correlation between vaping and psychiatric comorbidities, including generalized anxiety disorder, major depressive disorder, and attention-deficit/hyperactivity disorder (ADHD), suggesting a pattern of self-medication or shared neurobiological vulnerabilities. The epidemiology of EVALI is particularly distinct; during the outbreak peak in 2019 and 2020, thousands of hospitalizations and dozens of deaths were recorded, primarily affecting young males with a median age of 24 who reported using informal or illicitly sourced THC-containing vape cartridges.

\section*{Aetiology and Pathophysiology}
The pathological consequences of vaping arise from the inhalation of complex chemical mixtures and their interactions with the pulmonary parenchyma, bronchial airways, and systemic circulation. The aetiological factors and their corresponding pathophysiological mechanisms include:
\begin{itemize}
    \item \textbf{Nicotine Toxicity and Addiction}: Nicotine is highly lipid-soluble and rapidly absorbed through the pulmonary vasculature, reaching the brain within seconds. It binds to nicotinic acetylcholine receptors (nAChRs) in the central nervous system, stimulating the release of dopamine, norepinephrine, and other neurotransmitters. In the developing adolescent brain, this leads to structural changes in the prefrontal cortex, impairing impulse control, executive function, and mood regulation, while establishing a powerful loop of physical dependence and withdrawal.
    \item \textbf{Aerosolized Excipients (Propylene Glycol and Vegetable Glycerin)}: While these substances are generally recognized as safe for oral ingestion, their inhalation represents a different toxicological profile. When exposed to the high heat of vaping coils, propylene glycol and vegetable glycerin undergo thermal degradation (pyrolysis), producing reactive carbonyl compounds including formaldehyde, acetaldehyde, and acrolein. Acrolein is a potent airway irritant that damages the mucosal lining, impairs ciliary beat frequency, and induces bronchial hyperresponsiveness.
    \item \textbf{Flavoring Chemical Compounds}: Vaping liquids contain numerous flavoring agents (e.g., diacetyl, acetoin, acetyl propionyl, benzaldehyde, cinnamaldehyde) to produce sweet, fruity, or minty profiles. Diacetyl and acetyl propionyl are known to cause severe inflammatory and fibrotic changes in the distal airways, leading to bronchiolitis obliterans ("popcorn lung"). These flavorings induce cellular oxidative stress, trigger the release of pro-inflammatory cytokines (such as interleukin-8), and disrupt mitochondrial function in bronchial epithelial cells.
    \item \textbf{Heavy Metal Contamination}: The heating coils in vaping devices, typically composed of nichrome, kanthal, or stainless steel, leach heavy metals into the e-liquid during operation. Inhaled aerosols contain detectable levels of nickel, chromium, lead, manganese, and copper. These metals accumulate in pulmonary tissues, generating reactive oxygen species (ROS) that cause oxidative DNA damage, lipid peroxidation of cellular membranes, and chronic parenchymal inflammation.
    \item \textbf{Vitamin E Acetate}: Historically used as a thickening agent in illicit THC-containing vaping products, Vitamin E acetate is the primary driver of EVALI. When inhaled, its highly lipophilic structure allows it to insert into the alveolar surfactant monolayer. This disrupts the normal surface tension-reducing properties of the surfactant, leading to alveolar instability, collapse (atelectasis), and direct chemical pneumonitis. The physical presence of exogenous lipids within the alveoli triggers an influx of alveolar macrophages, which engulf the lipids, forming characteristic lipid-laden macrophages.
    \item \textbf{Tetrahydrocannabinol (THC) and Synthetic Cannabinoids}: The inhalation of vaporized cannabinoids can induce acute vasoactive and immunological effects. Synthetic cannabinoids, which are highly potent full agonists of cannabinoid receptors (CB1 and CB2), can cause severe cardiovascular instability, acute psychosis, and direct alveolar-capillary membrane disruption, leading to pulmonary edema and diffuse alveolar hemorrhage.
\end{itemize}

\section*{Clinical Features}
The clinical presentation of vaping-related disorders ranges from chronic, low-grade mucosal irritation to acute, life-threatening respiratory failure. Practitioners must recognize a wide spectrum of signs and symptoms:
\begin{itemize}
    \item \textbf{Respiratory Signs and Symptoms}: The most prominent features are progressive dyspnoea (shortness of breath), which may initially occur only on exertion but can progress to dyspnoea at rest. Patients frequently present with a dry, non-productive cough, pleuritic chest pain or tightness, tachypnoea (rapid breathing), and hypoxia (low blood oxygen saturation). In severe cases of diffuse alveolar hemorrhage or necrotizing pneumonitis, hemoptysis (coughing up blood) may be observed. Physical exam may reveal bibasilar crackles or wheezing, although auscultation can sometimes remain deceptively normal even in the presence of significant lung injury.
    \item \textbf{Gastrointestinal Symptoms}: A unique clinical feature of acute vaping-related lung injuries like EVALI is the high prevalence of gastrointestinal distress. Patients frequently experience severe nausea, persistent vomiting, abdominal pain, and diarrhea. These symptoms often precede the onset of respiratory complaints by several days or weeks, frequently leading to initial misdiagnoses of gastroenteritis or acute abdomen.
    \item \textbf{Constitutional and Systemic Features}: Patients often display systemic inflammatory responses, including subjective fever, chills, rigors, profound fatigue, malaise, and night sweats. Significant, unexplained weight loss is common, reflecting a chronic hyperinflammatory state and poor oral intake due to gastrointestinal distress.
    \item \textbf{Neurological and Neurocognitive Features}: Due to the high nicotine content in modern pod-based systems, acute nicotine toxicity can present with dizziness, headache, tremors, and palpitations. Conversely, nicotine withdrawal presents as extreme irritability, anxiety, difficulty concentrating, sleep disturbances, and intense drug cravings.
    \item \textbf{Cardiovascular Features}: Acute inhalation of nicotine-containing aerosol leads to immediate sympathetic nervous system activation, manifested as sinus tachycardia, elevated systolic and diastolic blood pressure, and acute peripheral vasoconstriction.
\end{itemize}

\section*{Differential Diagnosis}
Because the clinical features of vaping-induced lung injuries and nicotine toxicity are non-specific, clinicians must maintain a broad differential diagnosis. The alternative diagnoses to consider include:
\begin{itemize}
    \item \textbf{Community-Acquired Pneumonia (CAP)}: Bacterial or atypical infections (e.g., \textit{Streptococcus pneumoniae}, \textit{Mycoplasma pneumoniae}) and viral infections (such as Influenza, Respiratory Syncytial Virus [RSV], and COVID-19) present with overlapping symptoms of fever, cough, and dyspnoea. Differentiation relies on microbiologic cultures, PCR viral panels, and radiological findings showing focal lobar consolidation rather than the diffuse, ground-glass opacities with subpleural sparing typical of EVALI.
    \item \textbf{Asthma and COPD Exacerbation}: Exacerbations of pre-existing airway diseases can be triggered by vaping-induced bronchoconstriction. A history of childhood asthma, atopic features, prior spirometry showing reversible airway obstruction, and rapid clinical response to bronchodilators help distinguish these from primary vaping-induced parenchymal injury.
    \item \textbf{Hypersensitivity Pneumonitis (HP)}: This is an immunologically mediated inflammatory lung disease resulting from inhalation of environmental antigens. HP can mimic vaping injuries radiologically and clinically. A detailed occupational and environmental exposure history, specific serum IgG antibody testing, and the presence of poorly formed non-necrotizing granulomas on lung biopsy support a diagnosis of HP.
    \item \textbf{Cryptogenic Organizing Pneumonia (COP)}: COP is an idiopathic interstitial lung disease presenting with flu-like symptoms and patchy, bilateral consolidations. It is distinguished primarily by the absence of a history of toxic inhalational exposures and a typical lack of response to simple withdrawal of the exposure alone, requiring formal exclusion of other connective tissue and infectious etiologies.
    \item \textbf{Acute Respiratory Distress Syndrome (ARDS)}: ARDS represents a severe, diffuse inflammatory injury to the alveolar-capillary barrier. While severe EVALI can progress to ARDS, classic ARDS is typically precipitated by a clear systemic insult such as sepsis, severe trauma, aspiration of gastric contents, or pancreatitis.
    \item \textbf{Pneumocystis jirovecii Pneumonia (PJP)}: An opportunistic fungal infection that presents with progressive dyspnoea, dry cough, and diffuse ground-glass opacities, closely resembling vaping-induced lung injury. PJP is differentiated by the patient's immunocompromised status (e.g., advanced HIV/AIDS, immunosuppressive therapy) and positive induced sputum or BAL silver stains/PCR.
\end{itemize}

\section*{When to Seek Medical Care}
Vaping-related health issues can escalate rapidly from mild irritation to life-threatening respiratory failure. Individuals who vape should seek immediate medical evaluation if they develop any of the following warning signs:
\begin{itemize}
    \item \textbf{Emergency Care Indicators}: Immediate emergency services must be contacted if the individual experiences severe, sudden-onset shortness of breath or difficulty breathing, persistent chest pain or pressure, a bluish discoloration of the lips, face, or nail beds (cyanosis), confusion, disorientation, inability to stay awake, or sudden loss of consciousness. For life-threatening symptoms, emergency services should be contacted immediately by dialing:
    \begin{itemize}
        \item \textbf{local emergency services} in many countries.
        \item \textbf{Nine One One (911)} in the United States.
        \item \textbf{Triple Nine (999)} in the United Kingdom.
        \item \textbf{One One Two (112)} as the standard European emergency number.
    \end{itemize}
    \item \textbf{Routine Medical Consultation}: A doctor or general practitioner should be contacted promptly if the individual experiences a persistent, unexplained cough, gradual increase in shortness of breath during daily activities, recurrent episodes of nausea, vomiting, or diarrhea, unexplained weight loss, night sweats, or ongoing signs of nicotine withdrawal that interfere with daily functioning.
\end{itemize}

\section*{Diagnosis and Investigations}
Diagnosing vaping-related injuries requires a high index of clinical suspicion, a meticulous exposure history, and a structured diagnostic workup. Because there is no single diagnostic test for vaping-induced lung injury, the diagnosis is largely one of exclusion:
\begin{itemize}
    \item \textbf{Exposure History}: The clinician must obtain a detailed history of vaping habits, including the specific devices used, types of e-liquids (nicotine concentration, THC, or other cannabinoids), flavorings, frequency of use, and whether the products were obtained from licensed dispensaries or informal, illicit sources.
    \item \textbf{Physical Examination}: A comprehensive physical exam should focus on vital signs. Particular attention must be paid to pulse oximetry (both at rest and during ambulation to detect exertional desaturation), respiratory rate (looking for tachypnoea), heart rate (tachycardia), and temperature (fever). Chest auscultation should be performed to detect crackles or wheezes.
    \item \textbf{Laboratory Investigations}: Initial blood tests should include a complete blood count (CBC), which frequently reveals leukocytosis (elevated white blood cell count) with a neutrophil predominance. Inflammatory markers, such as C-reactive protein (CRP) and erythrocyte sedimentation rate (ESR), are typically markedly elevated. Renal function panel, liver function tests, and coagulation profiles are necessary to assess for multi-organ involvement. A urine toxicology screen is essential to detect the presence of cannabinoids or other illicit substances.
    \item \textbf{Radiological Imaging}: A chest radiograph (CXR) is the initial imaging modality, often demonstrating bilateral, diffuse, or patchy interstitial and alveolar infiltrates. However, a high-resolution chest computed tomography (HRCT) scan is the gold standard for characterizing vaping-induced lung injury. The classic CT finding in EVALI is diffuse, bilateral ground-glass opacities, frequently displaying subpleural sparing (a rim of normal lung tissue immediately adjacent to the pleura) and occasional lobular consolidation.
    \item \textbf{Bronchoscopy and Bronchoalveolar Lavage (BAL)}: In severe or diagnostic-uncertainty cases, a bronchoscopy with BAL is performed. This procedure is critical to rule out atypical infections. The lavage fluid can be stained with Oil Red O to identify lipid-laden macrophages, which supports the diagnosis of lipoid pneumonia, a sub-type of vaping-induced injury. In suspected EVALI cases, BAL fluid can also be analyzed for the presence of vitamin E acetate.
\end{itemize}

\section*{Management}
The management of vaping-induced health conditions involves immediate cessation of vape product use, supportive medical therapy, pharmacological intervention, and structured long-term cessation support:
\begin{itemize}
    \item \textbf{Hospitalization and Supportive Care}: Patients presenting with hypoxia (oxygen saturation below 92\% to 95\% on room air) or respiratory distress should be admitted to the hospital. Supportive respiratory management includes supplemental oxygen via nasal cannula or high-flow nasal oxygen (HFNO). In severe cases of respiratory failure or ARDS, non-invasive positive pressure ventilation (NIPPV) or endotracheal intubation and mechanical ventilation may be required. Refractory cases may necessitate extracorporeal membrane oxygenation (ECMO).
    \item \textbf{Pharmacological Therapy}: Corticosteroids are the primary pharmacologic treatment for moderate-to-severe vaping-induced lung injury (EVALI), as they rapidly reduce pulmonary inflammation. A typical regimen involves intravenous methylprednisolone (e.g., 1 to 2 mg/kg/day) during the acute phase, followed by a gradual oral taper of prednisone (starting at 40 to 60 mg/day) over 2 to 6 weeks, tailored to clinical response and imaging resolution. Because clinical presentation overlaps with infectious pneumonia, empiric broad-spectrum antibiotics and/or antivirals (e.g., oseltamivir) are usually initiated and subsequently discontinued once infectious cultures and PCR results return negative.
    \item \textbf{Treatment of Nicotine Dependence}: To manage severe nicotine withdrawal during acute illness and hospitalization, clinicians should offer nicotine replacement therapy (NRT) in the form of transdermal patches, nicotine gum, lozenges, or inhalers. Long-term cessation aids such as varenicline or bupropion can be considered post-discharge, provided there are no contraindications.
    \item \textbf{Behavioral and Psychological Support}: Enrollment in structured smoking and vaping cessation programs is vital. This includes cognitive behavioral therapy (CBT), motivational interviewing, and referral to digital support resources (e.g., Quitline services, SMS-based cessation programs). These programs help patients develop coping mechanisms for cravings, address underlying psychological distress, and prevent relapse.
    \item \textbf{Clinical Follow-up}: Patients discharged after a vaping-related lung injury require close outpatient monitoring. A follow-up visit should be scheduled within 1 to 2 weeks, including a physical exam, pulse oximetry, and spirometry (pulmonary function testing) to assess recovery of lung volumes. A repeat chest radiograph or CT scan should be performed within 2 to 4 weeks to document complete resolution of infiltrates and ensure no development of organizing pneumonia or progressive fibrosis.
\end{itemize}

\section*{Complications}
Untreated vaping-related disorders, or continued vaping after an injury, can lead to severe, irreversible systemic and pulmonary complications:
\begin{itemize}
    \item \textbf{Acute Respiratory Distress Syndrome (ARDS)}: Rapid progression of vaping-induced inflammation can lead to diffuse alveolar damage, severe refractory hypoxemia, multi-organ failure, and death.
    \item \textbf{Spontaneous Pneumothorax and Pneumomediastinum}: The high negative intrathoracic pressures generated during deep, forceful inhalation of vaping aerosols, combined with drug-induced alveolar wall weakening, can cause alveolar rupture. This leads to air escaping into the pleural cavity (pneumothorax) or the mediastinum (pneumomediastinum), requiring urgent chest tube insertion.
    \item \textbf{Popcorn Lung (Bronchiolitis Obliterans)}: Inhalation of flavoring chemicals like diacetyl causes severe, irreversible scarring and narrowing of the bronchioles. This results in chronic, fixed airway obstruction, characterized by progressive dyspnoea, wheezing, and a dry cough that does not respond to standard asthma medications.
    \item \textbf{Chronic Pulmonary Fibrosis}: Repetitive chemical insults to the lung parenchyma can trigger abnormal wound-healing pathways, leading to permanent scarring (fibrosis) of the interstitial lung tissue, chronic oxygen dependence, and severe restriction in functional capacity.
    \item \textbf{Cardiovascular Disease}: Chronic nicotine inhalation and exposure to fine particulate matter induce persistent endothelial dysfunction, increased arterial stiffness, and chronic systemic inflammation. This accelerates atherogenesis and significantly elevates the long-term risk of myocardial infarction, stroke, and peripheral arterial disease.
    \item \textbf{Accidental Nicotine Poisoning}: Concentrated e-liquids present a major risk of accidental ingestion or dermal absorption, particularly in toddlers and young children. Acute nicotine poisoning can cause vomiting, cardiac arrhythmias, seizures, respiratory depression, and death.
\end{itemize}

\section*{Prognosis and Outcomes}
The prognosis for individuals experiencing vaping-related health conditions is highly dependent on the severity of the initial presentation, the promptness of medical intervention, and the individual's ability to maintain complete abstinence from vaping. For patients diagnosed with acute lung injuries such as EVALI, the short-term prognosis is generally favorable if treatment is initiated early. More than 90\% of patients hospitalized with EVALI survive to discharge, with the vast majority showing rapid clinical and radiological improvement within days of starting corticosteroid therapy. However, patients who require mechanical ventilation or ECMO face a significantly higher risk of mortality and prolonged intensive care unit stays.

Long-term outcomes remain a subject of active medical investigation. Many patients who recover from severe vaping-induced lung injuries continue to exhibit mild-to-moderate pulmonary dysfunction, reduced exercise tolerance, and persistent respiratory symptoms (such as exertional dyspnoea or cough) for several months post-discharge. Radiologically, some individuals display residual scarring or subclinical interstitial changes, raising concerns about the potential for early-onset chronic obstructive lung disease. Crucially, the risk of recurrence of acute lung injury is high among those who resume vaping after discharge. The prognosis for nicotine addiction acquired through vaping is variable; while many succeed in quitting with combined pharmacotherapy and behavioral support, the high nicotine delivery of modern devices makes relapse common, exposing individuals to sustained cardiovascular and carcinogenic risks over their lifespan.

\section\*{Prevention and Self-Care}
Preventing vaping-related illnesses and maintaining long-term respiratory health requires comprehensive public health measures, individual risk mitigation, and proactive self-care strategies:
\begin{itemize}
    \item \textbf{Absolute Avoidance and Cessation}: The primary preventive measure is the complete avoidance of all vaping products, e-cigarettes, traditional cigarettes, and heated tobacco products. Non-users, particularly adolescents, pregnant women, and young adults, should never initiate vaping.
    \item \textbf{Avoidance of Illicit and Informal Products}: Under no circumstances should individuals purchase vaping products, e-liquids, or cartridges from informal sources, street vendors, online illicit markets, or modify commercial products to add substances like THC or synthetic oils.
    \item \textbf{Home Environmental Control}: Individuals should ensure their living and working environments are completely smoke-free and vape-free. Secondhand exposure to vaping aerosol contains harmful particulates and chemicals that can irritate sensitive airways, particularly in children and individuals with asthma.
    \item \textbf{Self-Monitoring and Early Detection}: Recovering patients should perform daily self-monitoring. This includes tracking respiratory symptoms (cough, shortness of breath), checking body temperature, and using a home pulse oximeter to monitor resting and exertional oxygen saturation. Any drop in oxygen saturation below 95\% or a return of gastrointestinal or respiratory symptoms should trigger immediate medical re-evaluation.
    \item \textbf{Healthy Pulmonary Habits}: To aid lung recovery and maintain overall health, individuals should engage in regular, low-impact cardiovascular exercise as tolerated (e.g., walking, swimming) to improve ventilatory efficiency. Practicing deep breathing exercises and ensuring adequate hydration helps thin bronchial secretions and supports the mucociliary escalator.
    \item \textbf{Utilization of Cessation Networks}: Individuals trying to quit should actively engage with professional support networks. This includes registering with national Quitline services, utilizing evidence-based mobile applications for behavioral tracking, and consulting healthcare professionals to devise a personalized cessation plan.
\end{itemize}

\section*{Sources and Further Reading}
For authoritative information, clinical guidelines, and support resources regarding vaping and cessation, consult the following organizations:
\begin{itemize}
    \item World Health Organization (WHO) --- Tobacco Free Initiative: \url{https://www.who.int/health-topics/tobacco}
    \item Centers for Disease Control and Prevention (CDC) --- E-cigarette and Vaping Information: \url{https://www.cdc.gov/tobacco/basic-information/e-cigarettes/index.html}
    \item Australian Government Department of Health and Aged Care --- Vaping and E-cigarettes: \url{https://www.health.gov.au/topics/vaping-and-e-cigarettes}
    \item National Institutes of Health (NIH) / National Institute on Drug Abuse (NIDA) --- Vaping Devices (E-Cigarettes) DrugFacts: \url{https://nida.nih.gov/publications/drugfacts/vaping-devices-e-cigarettes}
    \item Mayo Clinic --- Vaping: What You Need to Know: \url{https://www.mayoclinic.org/healthy-lifestyle/quit-smoking/in-depth/vaping/art-20468926}
    \item Asthma + Lung UK --- Vaping and Lung Health: \url{https://www.asthmaandlung.org.uk/living-with/vaping}
    \item U.S. Food and Drug Administration (FDA) --- Vaporizers, E-Cigarettes, and other Electronic Nicotine Delivery Systems (ENDS): \url{https://www.fda.gov/tobacco-products/products-ingredients-components/vaporizers-e-cigarettes-and-other-electronic-nicotine-delivery-systems-ends}
\end{itemize}